\chapter[2025 August]{August 2025}

\section[2025/08/14]{Thursday, 14 August 2025}
\label{sec:20250814}

There are two main modifications made in  order to verify this simple GAT is working. These changes
are made in v2

\begin{compactitem}
    \item The first is the dataset which allows for the filtering of data. A new parameter is added
    num\_people\_required which when set to $-1$ will return the dataset as normal but when set to a number
    it returns that many people; for now, more specifically, it returns that many full skeletons where all the 
    keypoints are visible.
    \item I moved all the hyperparameters to settings (one place) and added a method to write those to disk.
    This way when training (in the trainer) I can generate a guid (a unique reference) per run and when I save
    the results I just save the hyperparameters as well so I know what generated those results
\end{compactitem}

There are fairly few images (in the training set) that has two complete humans (100 images) so the dataset 
is fairly small.

\newcommand{\runid}{a104176f-94db-4437-8047-6b2a5c628ca2}

\subsection{Run: \texttt{\runid}}

This is the first run and some parameters I've chosen at random

\lstinputlisting[
    caption={Hyperparameters for run \texttt{\runid}},
    label={code:hyperparameters:\runid}
]{figures/\runid/hyperparameters.json}

\begin{figure}[H]
  \centering
  \includegraphics[width=0.48\linewidth]{figures/\runid/embeddings_epoch_25.png}
  \includegraphics[width=0.48\linewidth]{figures/\runid/embeddings_epoch_50.png}
  \caption{Embeddings at various epochs for run \texttt{\runid}.}
  \label{fig:embeddings:\runid}
\end{figure}

\begin{figure}[H]
  \centering
  \includegraphics[width=0.9\linewidth]{figures/\runid/training_summary.png}
  \caption{Training summary for run \texttt{\runid}.}
  \label{fig:summary:\runid}
\end{figure}

\paragraph{Training Loss Dynamics} 
The contrastive loss decreased rapidly
during the first 10 epochs (from a high initial value down to $13.80$) and
continued to drop sharply to $1.15$ by epoch $20$. By epoch $50$, the loss
plateaued at approximately $0.84$, indicating early convergence. This behaviour
suggests that the model quickly satisfied the contrastive margin constraints
and subsequently made only marginal improvements.

\paragraph{Positive and Negative Distances} 
The average positive-pair distance
(\textit{same person}) fell from $13.79$ at epoch $10$ to $0.40$ at epoch $50$.
This shows that embeddings for keypoints belonging to the same person were
pulled close together effectively.

The average negative-pair distance (\textit{different persons}) also fell
steeply, from $15.44$ at epoch $10$ to $0.64$ at epoch $50$. Ideally, negative
distances should remain large or increase to ensure strong inter-person
separation. Here, the reduction in negative distances caused the final
separation 
\[ \Delta d = d_{\mathrm{neg}} - d_{\mathrm{pos}} \approx 0.24 \]
to be small, indicating poor discrimination between persons.

\paragraph{Clustering Behaviour}

PCA visualisation of embeddings at epochs $25$ and $50$ revealed that points
were concentrated along a near one-dimensional subspace, with the first
principal component explaining almost all variance and the second component
contributing very little. Low variance in higher principal components indicates
that the embeddings collapse into a narrow manifold, limiting their
separability. This explains the small separation $\Delta d$: while intra-person
distances are small, inter-person distances are not sufficiently larger,
leading to overlapping clusters.

\paragraph{Diagnosis}
The small final separation, rapid early convergence, and near-1D embedding collapse suggest that:
\begin{itemize}
    \item The contrastive \textbf{margin} ($1.0$) is too low, allowing early satisfaction of the loss without forcing sufficient inter-person separation.
    \item The model capacity (\texttt{hidden\_dim}=32, \texttt{output\_dim}=16, \texttt{num\_heads}=2) may be insufficient to extract richer discriminative features from the input.
    \item The dropout rate ($0.1$) may be too low to regularise the model, contributing to embedding collapse.
\end{itemize}

\renewcommand{\runid}{46cfff8e-3db4-41ef-a0e8-925330f73c71}

\subsection{Run: \texttt{\runid}}

While it's probably better to make small incremental changes a number of drastic changes were made for this run

\lstinputlisting[
    caption={Hyperparameters for run \texttt{\runid}},
    label={code:hyperparameters:\runid}
]{figures/\runid/hyperparameters.json}

\begin{figure}[H]
  \centering
  \includegraphics[width=0.48\linewidth]{figures/\runid/embeddings_epoch_25.png}
  \includegraphics[width=0.48\linewidth]{figures/\runid/embeddings_epoch_50.png}
  \includegraphics[width=0.48\linewidth]{figures/\runid/embeddings_epoch_75.png}
  \caption{Embeddings at various epochs for run \texttt{\runid}.}
  \label{fig:embeddings:\runid}
\end{figure}

\begin{figure}[H]
  \centering
  \includegraphics[width=0.9\linewidth]{figures/\runid/training_summary.png}
  \caption{Training summary for run \texttt{\runid}.}
  \label{fig:summary:\runid}
\end{figure}

\paragraph{Training Loss Dynamics}

The contrastive loss began extremely high ($\approx 429$) and dropped sharply
within the first 10 epochs to $28.90$. By epoch $20$, the loss was already down
to $2.31$, after which it decreased only slightly, plateauing near $1.90$ from
epoch $40$ onward. This pattern indicates rapid early convergence: during the
initial epochs, positive pairs were pulled close together and most negative
pairs were pushed apart to at least the specified margin ($m = 2.0$). Once a
negative pair's distance exceeds the margin, its contribution to the loss
becomes zero, meaning it no longer exerts a repulsive force. As a result, after
the first $\sim 20$ epochs, the majority of pairs were already within their
target distances, leaving the model with little gradient signal to further
increase inter-person separation, causing the loss to plateau.

\paragraph{Positive and Negative Distances}

The average positive-pair distance (\textit{same person}) decreased quickly
from $28.72$ at epoch $10$ to $1.01$ by epoch $20$ and then stabilised around
$0.58$ from epoch $40$ onwards.

The average negative-pair distance (\textit{different persons}) also dropped
rapidly from $29.62$ at epoch $10$ to $1.12$ by epoch $20$, then stabilised
around $0.68$ by the end of training.

The final separation was:
\[
\Delta d = d_{\mathrm{neg}} - d_{\mathrm{pos}} \approx 0.11
\]

This extremely small separation shows that negatives remain nearly as close as
positives in embedding space, which significantly limits clustering
performance.

\paragraph{Clustering Behaviour}

PCA plots at epochs $25$, $50$, and $75$ reveal a collapse of embedding
variance into an almost single-dimensional structure. At epoch $25$, the first
principal component explained $72.47\%$ of variance, with the second capturing
$27.03\%$, showing some spread. By epoch $50$, PC1 explained essentially all
variance ($100\%$), with PC2 contributing $0\%$, indicating a complete collapse
to a line. Epoch $75$ showed a similar trend ($99.55\%$ variance in PC1,
$0.42\%$ in PC2).
Low variance in higher principal components means that embeddings are being
projected into a very narrow manifold, leaving little capacity for separating
individuals in multiple directions of the embedding space.

\paragraph{Diagnosis}
While the increased margin ($2.0$) and model capacity (64 hidden units, 32-dimensional embeddings, 4 attention heads) prevented immediate trivial collapse, the training behaviour indicates:
\begin{itemize}
    \item Negatives are still being pulled too close, suggesting that the loss weighting or formulation may not sufficiently penalise small negative distances.
    \item The network may be overfitting early to a degenerate embedding pattern where both positives and negatives collapse together, satisfying the margin in minimal dimensions.
    \item Even with more capacity, the embeddings are not making use of the additional representational space.
\end{itemize}

\subsection{PCA aside}

The PCA is a way to re-express the embeddings into a new (simpler coordinate system)

\begin{compactitem}
\item The first axis (PC1) captures as much variance in the data as possible 
\item The second axis (PC2) captures the largest amount of remaining variation, but is orthogonal (at a right angle) to PC1
\item The third axis (PC3) captures the next largest amount of remaining variance, and is orthogonal to both PC1 and PC2
\item etc
\end{compactitem}

This matters because
\begin{compactitem}
\item The eigenvalues in PCA tell you how much of the total variation in the data is captured by each principle component
\item If PC1 is 99\% of the variance for example (like in some of the runs) it means points are almost entirely spread out along one direction in space with little spread in other directions
\end{compactitem}

Good embeddings for clustering should spread out in multiple independent directions

\renewcommand{\runid}{6f93a28e-9121-40d2-bdcb-9584ffa73d66}

\subsection{Run: \texttt{\runid}}

\lstinputlisting[
    caption={Hyperparameters for run \texttt{\runid}},
    label={code:hyperparameters:\runid}
]{figures/\runid/hyperparameters.json}

\begin{figure}[H]
  \centering
  \includegraphics[width=0.48\linewidth]{figures/\runid/embeddings_epoch_25.png}
  \includegraphics[width=0.48\linewidth]{figures/\runid/embeddings_epoch_50.png}
  \includegraphics[width=0.48\linewidth]{figures/\runid/embeddings_epoch_75.png}
  \caption{Embeddings at various epochs for run \texttt{\runid}.}
  \label{fig:embeddings:\runid}
\end{figure}

\begin{figure}[H]
  \centering
  \includegraphics[width=0.9\linewidth]{figures/\runid/training_summary.png}
  \caption{Training summary for run \texttt{\runid}.}
  \label{fig:summary:\runid}
\end{figure}

\paragraph{Training Loss Dynamics} 

The contrastive loss started very high ($\approx 100.47$ at epoch~10) and
dropped sharply within the first 20~epochs to $10.14$. It continued decreasing
at a slower rate, reaching $5.85$ at epoch~30 and plateauing near $4.97$ from
epoch~60 onward. This plateau indicates that the model reached a state where
most negative pairs had distances exceeding the margin ($m = 5.0$), resulting
in negligible loss contributions from them and limiting further separation.

\paragraph{Positive and Negative Distances} 

The average positive-pair distance decreased steadily from $99.58$ at epoch~10
to $0.65$ at epoch~70, showing strong intra-person compactness. The average
negative-pair distance also fell dramatically, from $100.94$ at epoch~10 to
$0.71$ at epoch~70. This parallel decrease indicates that while positives were
pulled together effectively, negatives were not kept far apart, leading to a
final separation

\[
\Delta d = d_{\mathrm{neg}} - d_{\mathrm{pos}} \approx 0.0858,
\]

which is insufficient for clear inter-person discrimination.

\paragraph{Clustering Behaviour} 

PCA visualisations show progressive variance collapse into a single principal
component. At epoch~25, the first two components explained $48.93\%$ and
$36.26\%$ of the variance, suggesting a relatively balanced spread in two
dimensions. By epoch~50, the second component’s variance dropped to $18.08\%$,
and by epoch~75, the first component explained $99.98\%$ of the variance, with
the second contributing almost none. This near one-dimensional collapse means
that embeddings are aligned along a single axis, severely limiting the ability
to form well-separated clusters in the embedding space.

\paragraph{Diagnosis} 

The rapid drop in both positive and negative distances, coupled with PCA
collapse, suggests that the model learned to indiscriminately compress all
embeddings along one dominant axis. Contributing factors likely include:

\begin{itemize}
    \item Margin $m = 5.0$ being reached quickly for negatives, halting repulsion.
    \item A learning rate of $3 \times 10^{-4}$ allowing fast early compression without preserving inter-class variance.
    \item Dropout $0.4$ possibly regularising too strongly, encouraging low-dimensional representations.
\end{itemize}

\renewcommand{\runid}{ae3b40be-6ed7-4fd4-862e-7f1911bf05cb}

\subsection{Run: \texttt{\runid}}

\lstinputlisting[
    caption={Hyperparameters for run \texttt{\runid}},
    label={code:hyperparameters:\runid}
]{figures/\runid/hyperparameters.json}

\begin{figure}[H]
  \centering
  \includegraphics[width=0.48\linewidth]{figures/\runid/embeddings_epoch_25.png}
  \includegraphics[width=0.48\linewidth]{figures/\runid/embeddings_epoch_50.png}
  \includegraphics[width=0.48\linewidth]{figures/\runid/embeddings_epoch_75.png}
  \caption{Embeddings at various epochs for run \texttt{\runid}.}
  \label{fig:embeddings:\runid}
\end{figure}

\begin{figure}[H]
  \centering
  \includegraphics[width=0.9\linewidth]{figures/\runid/training_summary.png}
  \caption{Training summary for run \texttt{\runid}.}
  \label{fig:summary:\runid}
\end{figure}

\paragraph{Run Summary}

This experiment increased the model capacity and margin relative to earlier
runs:

\begin{itemize}
    \item \texttt{hidden\_dim} increased from $32$ to $128$, and \texttt{num\_heads} from $2$ to $8$, to enable richer multi-head attention patterns.
    \item The contrastive loss \texttt{margin} was set to $10.0$ (previously $1.0$ or $3.0$), forcing the model to push negative pairs much further apart in embedding space.
    \item Dropout increased to $0.4$ to counter potential overfitting from the larger model.
\end{itemize}

The training loss dropped rapidly from $\approx 250$ to $\approx 9.3$ by epoch
$20$, then plateaued. Final metrics were:

\[
d_{\mathrm{pos}} = 2.97, \quad d_{\mathrm{neg}} = 3.95, \quad \Delta d = 0.99
\]

showing a large improvement in separation compared to earlier runs.

PCA visualisations revealed that although the first principal component still
dominates variance ($\approx 95$--$99\%$), the increased margin prevented
complete embedding collapse: clusters for different persons were more clearly
separated along PC1, and $\Delta d$ increased by an order of magnitude relative
to the baseline. Also interestingly the neg distance was slowly climbing

\section[2025/08/15]{Friday, 15 August 2025}
\label{sec:20250815}

This is obviously not working as expected. There are two alternative approaches to try which might improve 
inter-person separation

\subsection{Option A: L2 Distance with Squared-Hinge Loss} In the current
setup, the contrastive loss for negative pairs uses a standard hinge
formulation:

\[
\ell_{\text{neg}} = \max(0, m - d),
\]

where \(m\) is the contrastive margin and \(d\) is the Euclidean distance
between embeddings. Once \(d \geq m\), the loss becomes zero, and the model
stops pushing the negative pair apart. In prior runs, this caused many
negatives to reach the margin quickly, removing their contribution to the
gradient and allowing their distances to drift back down, reducing inter-person
separation.

The squared-hinge modification replaces the loss term with:

\[
\ell_{\text{neg}} = \max(0, m - d)^2,
\]

which has two key effects:

\begin{enumerate}
    \item \textbf{Stronger repulsion near the margin:} Squaring the term increases the gradient when \(d\) is close to \(m\), encouraging the model to continue pushing negatives apart instead of settling just at the margin.
    \item \textbf{Smoother gradient decay:} The gradient reduces more gradually as \(d\) approaches \(m\), making the optimisation less abrupt and potentially more stable.
\end{enumerate}

Here the margin \(m\) will be increased (from $2.0$ to $3.0$ or $4$ etc), the
hidden and output dimensions are enlarged to increase model capacity, and the
negative-pair loss term is up-weighted (\texttt{neg\_weight} = 2.0) so that
repulsion receives greater emphasis during training. Together, these changes
are intended to prevent early “margin satisfaction,” maintain pressure on
separating negatives, and reduce the embedding collapse into a near
one-dimensional structure observed in PCA.


\subsection{Option B: Cosine Margin Loss with Normalised Embeddings}

In this alternative approach, the embedding vectors produced by the GAT are
first $L_{2}$-normalised:

\[
\mathbf{z} \leftarrow \frac{\mathbf{z}}{\lVert \mathbf{z} \rVert_2},
\]

ensuring that all embeddings lie on the unit hypersphere. This normalisation
fixes the scale of the vectors, making their relative position depend only on
\emph{direction} (angle) rather than magnitude. In previous runs, embeddings
often collapsed in scale and aligned along a single dominant axis, as seen in
PCA variance results. By constraining all vectors to have the same length, the
model is prevented from trivially shrinking embedding magnitudes to reduce
distances, forcing it to learn meaningful angular separations.

The loss function is replaced with a cosine-margin formulation:

\[
\ell_{\text{pos}} = 1 - \cos(\theta_{ij}), \quad
\ell_{\text{neg}} = \max\big(0, \cos(\theta_{ij}) - (1 - m)\big),
\]

where $\theta_{ij}$ is the angle between two embeddings and $m$ is the cosine
margin. Positive pairs minimise their angular distance (maximising cosine
similarity), while negative pairs are penalised if their cosine similarity
exceeds $1 - m$. For $m = 0.4$, negatives are encouraged to have an angular
separation of at least $\arccos(0.6) \approx 53^\circ$.

This angular-based formulation mitigates the collapse problem in two ways:
\begin{enumerate}
    \item \textbf{Magnitude invariance:} Because vectors are normalised, the network cannot exploit magnitude shrinkage to satisfy the loss, as occurred in earlier L2-based runs.
    \item \textbf{Direct angular separation:} Enforcing a fixed angular margin ensures negatives remain well-separated in direction, even if their Euclidean distances would otherwise converge.
\end{enumerate}

Model capacity is also increased (\texttt{hidden\_dim} = 96, \texttt{output\_dim} = 64, \texttt{num\_heads} = 8) to provide more representational diversity, and dropout is set to $0.25$ to balance regularisation with the need to capture multi-directional structure in the embedding space.

\subsection{Simplification}

For a test use just x,y for node features (same metric) to see if the embeddings move apart

\section[2025/08/14]{Thursday, 28 August 2025}
\label{sec:20250828}

Doing the simplification where only x,y coordinates are used for node features (same metric) to see if the embeddings move apart.

\renewcommand{\runid}{25be564b-ee51-407d-a663-a7edbe0f51f6}

\subsection{Run: \texttt{\runid}}

\lstinputlisting[
    caption={Hyperparameters for run \texttt{\runid}},
    label={code:hyperparameters:\runid}
]{figures/\runid/hyperparameters.json}

\begin{figure}[H]
  \centering
  \includegraphics[width=0.48\linewidth]{figures/\runid/embeddings_epoch_25.png}
  \includegraphics[width=0.48\linewidth]{figures/\runid/embeddings_epoch_50.png}
  \caption{Embeddings at various epochs for run \texttt{\runid}.}
  \label{fig:embeddings:\runid}
\end{figure}

\begin{figure}[H]
  \centering
  \includegraphics[width=0.9\linewidth]{figures/\runid/training_summary.png}
  \caption{Training summary for run \texttt{\runid}.}
  \label{fig:summary:\runid}
\end{figure}

\paragraph{Metrics:}

Loss steadily decreases from 8.55 to 6.72 (model is learning). 

Positive distances decrease 5.01 to 3.28 which means the same person keypoints are getting closer in embedding space

Negative distances increase 7.19 to 7.94 different person keypoints are being pushed further apart

\paragraph{PCA:} The PCA output shows:

\begin{compactitem}
  \item PC1 explains around 99.9\% of the variance.
  \item PC2 around 0.01\%
\end{compactitem}

This means that almost all structure of the embedding is in one dimension

(this can be seen on the 50 epoch embeddings where the cloud looks like a line stretched along PC1 with little variance in PC2)

Since PC1 dominates the model may effectively be collapsing embeddings into a 1D manifold: all useful information in one dimension

This can happen for the following reasons:

\begin{compactitem}
  \item small feature input (just x and y) model might only find one discriminative axis (like left-rigth spread of the keypoints)
  \item contrastive loss pushes/pulls distances but doesn't enforce mutli-dimensional use of embedding space 
  \item PCA reveals that almost all vairance is captured in one linear axis
\end{compactitem}


\paragraph{Next steps:} To avoid "all in a line" embeddings

\begin{compactitem}
  \item Increase feature set agtain (reintroduce depth, confidence, joint type etc)
  \item Normalize embeddings (L2 Norm) before contrastive loss - might avoid collapse 
  \item Increase embeddings dimension so it has the capacity to spread clusters into more directions
  \item try different loss 
\end{compactitem}


\renewcommand{\runid}{0211e02a-4fc1-414f-b7c9-9a8135507ca1}

\subsection{Run: \texttt{\runid}}

\lstinputlisting[
    caption={Hyperparameters for run \texttt{\runid}},
    label={code:hyperparameters:\runid}
]{figures/\runid/hyperparameters.json}

\begin{figure}[H]
  \centering
  \includegraphics[width=0.48\linewidth]{figures/\runid/embeddings_epoch_25.png}
  \includegraphics[width=0.48\linewidth]{figures/\runid/embeddings_epoch_50.png}
  \includegraphics[width=0.48\linewidth]{figures/\runid/embeddings_epoch_75.png}
  \includegraphics[width=0.48\linewidth]{figures/\runid/embeddings_epoch_100.png}
  \caption{Embeddings at various epochs for run \texttt{\runid}.}
  \label{fig:embeddings:\runid}
\end{figure}

\begin{figure}[H]
  \centering
  \includegraphics[width=0.9\linewidth]{figures/\runid/training_summary.png}
  \caption{Training summary for run \texttt{\runid}.}
  \label{fig:summary:\runid}
\end{figure}

similar to the last results.

loss decreases from 9.10 to 6.72, pos distance decreases from 6.23 to 3.45 and negative distance increases from 7.99 to 8.53

The PCA follows a similar trend.

\renewcommand{\runid}{d87e9c71-c64b-40d1-9689-fd0c7be9b128}

\subsection{Run: \texttt{\runid}}

This run adds back depth information to the keypoints to see if that makes a difference

\lstinputlisting[
    caption={Hyperparameters for run \texttt{\runid}},
    label={code:hyperparameters:\runid}
]{figures/\runid/hyperparameters.json}

\begin{figure}[H]
  \centering
  \includegraphics[width=0.48\linewidth]{figures/\runid/embeddings_epoch_25.png}
  \includegraphics[width=0.48\linewidth]{figures/\runid/embeddings_epoch_50.png}
  \includegraphics[width=0.48\linewidth]{figures/\runid/embeddings_epoch_75.png}
  \includegraphics[width=0.48\linewidth]{figures/\runid/embeddings_epoch_100.png}
  \caption{Embeddings at various epochs for run \texttt{\runid}.}
  \label{fig:embeddings:\runid}
\end{figure}

\begin{figure}[H]
  \centering
  \includegraphics[width=0.9\linewidth]{figures/\runid/training_summary.png}
  \caption{Training summary for run \texttt{\runid}.}
  \label{fig:summary:\runid}
\end{figure}

Pretty much the same results again.

\renewcommand{\runid}{e7d93493-1252-4117-8060-e47dbedf497e}

\subsection{Run: \texttt{\runid}}

This time using L2 Norm

\lstinputlisting[
    caption={Hyperparameters for run \texttt{\runid}},
    label={code:hyperparameters:\runid}
]{figures/\runid/hyperparameters.json}

\begin{figure}[H]
  \centering
  \includegraphics[width=0.48\linewidth]{figures/\runid/embeddings_epoch_25.png}
  \includegraphics[width=0.48\linewidth]{figures/\runid/embeddings_epoch_50.png}
  \includegraphics[width=0.48\linewidth]{figures/\runid/embeddings_epoch_75.png}
  \includegraphics[width=0.48\linewidth]{figures/\runid/embeddings_epoch_100.png}
  \caption{Embeddings at various epochs for run \texttt{\runid}.}
  \label{fig:embeddings:\runid}
\end{figure}

\begin{figure}[H]
  \centering
  \includegraphics[width=0.9\linewidth]{figures/\runid/training_summary.png}
  \caption{Training summary for run \texttt{\runid}.}
  \label{fig:summary:\runid}
\end{figure}

This didn't work as expected you can see from the neg - pos distance is increasing not decreasing.
This is probably just because the margin was still set at 10 with unit-norm embeddings. L2-normalized
vectors, Euclidean distance is bounded to $[0,2]$

\renewcommand{\runid}{d6508c32-9d31-46ae-b8d5-940e9c5f5660}

\subsection{Run: \texttt{\runid}}

\lstinputlisting[
    caption={Hyperparameters for run \texttt{\runid}},
    label={code:hyperparameters:\runid}
]{figures/\runid/hyperparameters.json}

\begin{figure}[H]
  \centering
  \includegraphics[width=0.48\linewidth]{figures/\runid/embeddings_epoch_25.png}
  \includegraphics[width=0.48\linewidth]{figures/\runid/embeddings_epoch_50.png}
  \includegraphics[width=0.48\linewidth]{figures/\runid/embeddings_epoch_75.png}
  \includegraphics[width=0.48\linewidth]{figures/\runid/embeddings_epoch_100.png}
  \caption{Embeddings at various epochs for run \texttt{\runid}.}
  \label{fig:embeddings:\runid}
\end{figure}

\begin{figure}[H]
  \centering
  \includegraphics[width=0.9\linewidth]{figures/\runid/training_summary.png}
  \caption{Training summary for run \texttt{\runid}.}
  \label{fig:summary:\runid}
\end{figure}

This still didn't work becaause distances are capped between 0 and 2 so setting margin=1 means that any negative pair closer than 1
contributes to losee and pairs father than 1 contribute to nothing

Guess you can tune the margin so set to 1.2 or 1.5 or something else so that negatives won't saturate too early but rather use
cosine similarity whicdh might work better then I'll try again

\renewcommand{\runid}{330bce9e-2578-485b-9ed5-d4ec9ff8072c}

\subsection{Run: \texttt{\runid}}

\lstinputlisting[
    caption={Hyperparameters for run \texttt{\runid}},
    label={code:hyperparameters:\runid}
]{figures/\runid/hyperparameters.json}

\begin{figure}[H]
  \centering
  \includegraphics[width=0.48\linewidth]{figures/\runid/embeddings_epoch_25.png}
  \includegraphics[width=0.48\linewidth]{figures/\runid/embeddings_epoch_50.png}
  \includegraphics[width=0.48\linewidth]{figures/\runid/embeddings_epoch_75.png}
  \includegraphics[width=0.48\linewidth]{figures/\runid/embeddings_epoch_100.png}
  \caption{Embeddings at various epochs for run \texttt{\runid}.}
  \label{fig:embeddings:\runid}
\end{figure}

\begin{figure}[H]
  \centering
  \includegraphics[width=0.9\linewidth]{figures/\runid/training_summary.png}
  \caption{Training summary for run \texttt{\runid}.}
  \label{fig:summary:\runid}
\end{figure}

This does also seem to be working. The training loss keeps decreasing (there are fluctuations). The embeddings similarities green (positive pair same person)
around 0.76 (want this close to 1 - where 1 is perfect alignment). The Red (negative pairs different people) around 0.38. There is a gap

The similarity seperation is now defined as pos - neg (in eucledian you have neg - pos)

\paragraph{Euclidean distance}: distances between embeddings

Make pos pairs (same person) small and neg pairs (different person) large (much bigger than margin)

\paragraph{Cosine similarity}: compute similarities using F.consine\_similarity

Range $[-1,1]$

positive pairs maximize similarity -> close to 1

negative pairs minimize similarity -> close to 0 or negative

so the metric becomes pos-neg (oppostive of Euclidean) and you want this to be large

looking at the PCA plots

\begin{compactitem}
  \item epoch 25: embeddings are fairly well-separated into two “blobs” (good clustering).
  \item epoch 50: points spread out more, not clean clusters
  \item epoch 75: separation comes back (two sides again).
  \item epoc 100: they spread into an arc shape, more variance across PC2.
\end{compactitem}

\begin{compactitem}
  \item PCA is linear and only keeps the top 2 components.
  \item If most of the separation is along 1 axis (say PC1), then PC2 may capture only tiny variance and look noisy. That's why the clouds look like a line or arc.
\end{compactitem}

Not sure what to think of this the logs seem fine pos sim = 0.76, neg sim = 0.38 and seperation = 0.37
\section[2025/08/14]{Friday, 29 August 2025}
\label{sec:20250829}

The results from yesterday somewhat prove that this method does work what I'm worried about now is that I process the dataset
meaning that there are (this is a setting thing) but 2 people where all points are visible and that by definition means that there is a naturual 
seperation.

Going to update this code v5 so that this isn't the case and it just selects 2 people even if their joints aren't all visible




Notes:
- use the temporal version of GUIDS (used in dbs) UILD 
- how do people generate skeltons and animate them 